% problem_id: S046-prop:trunc_of_prop_is_set
% source_id: S046 (HoTT Book)
% source_file: setmath.tex
% statement_locator: setmath.tex lines 537--540
% proof_locator: setmath.tex lines 542--612
% env: lem
% pairing: tight (verified)
% status: complete (statement + author proof verbatim + reason + locators)
%
\begin{lem}\label{prop:trunc_of_prop_is_set}
If $A$ is a mere proposition then its suspension $\susp(A)$ is a set,
and $A$ is equivalent to $\id[\susp(A)]{\north}{\south}$.
\end{lem}

\begin{proof}
To show that $\susp(A)$ is a set, we define a
family $P:\susp(A)\to\susp(A)\to\type$ with the
property that $P(x,y)$ is a mere proposition for each $x,y:\susp(A)$,
and which is equivalent to its identity type $\idtypevar{\susp(A)}$.
%
We make the following definitions:
\begin{align*}
P(\north,\north) & \defeq \unit &
P(\south,\north) & \defeq A\\
P(\north,\south) & \defeq A &
P(\south,\south) & \defeq \unit.
\end{align*}
We have to check that the definition preserves paths.
Given any $a : A$, there is a meridian $\merid(a) : \north = \south$,
so we should also have
%
\begin{equation*}
  P(\north, \north) = P(\north, \south) = P(\south, \north) = P(\south, \south).
\end{equation*}
%
But since $A$ is inhabited by $a$, it is equivalent to $\unit$, so we have
%
\begin{equation*}
  P(\north, \north) \eqvsym P(\north, \south) \eqvsym P(\south, \north) \eqvsym P(\south, \south).
\end{equation*}
%
The univalence axiom turns these into the desired equalities. Also, $P(x,y)$ is a mere
proposition for all $x, y : \susp(A)$, which is proved by induction on $x$ and $y$, and
using the fact that being a mere proposition is a mere proposition.

Note that $P$ is a reflexive relation.
Therefore we may apply \cref{thm:h-set-refrel-in-paths-sets}, so it suffices to
construct $\tau : \prd{x,y:\susp(A)}P(x,y)\to(x=y)$. We do this by a double induction.
When $x$ is $\north$, we define $\tau(\north)$ by
%
\begin{equation*}
  \tau(\north,\north,u) \defeq \refl{\north}
  \qquad\text{and}\qquad
  \tau(\north,\south,a) \defeq \merid(a).
\end{equation*}
%
If $A$ is inhabited by $a$ then $\merid(a) : \north = \south$ so we also need
\narrowequation{
  \trans{\merid(a)}{\tau(\north, \north)} = \tau(\north, \south).
}
This we get by function extensionality using the fact that, for all $x : A$,
%
\begin{multline*}
  \trans{\merid(a)}{\tau(\north,\north,x)} =
  \tau(\north,\north,x) \ct \opp{\merid(a)} \jdeq \\
  \refl{\north} \ct \merid(a) =
  \merid(a) =
  \merid(x) \jdeq
  \tau(\north, \south, x).
\end{multline*}
In a symmetric fashion we may define $\tau(\south)$ by
%
\begin{equation*}
  \tau(\south,\north, a) \defeq \opp{\merid(a)}
  \qquad\text{and}\qquad
  \tau(\south,\south, u) \defeq \refl{\south}.
\end{equation*}
%
To complete the construction of $\tau$, we need to check $\trans{\merid(a)}{\tau(\north)} = \tau(\south)$,
given any $a : A$. The verification proceeds much along the same lines by induction on the
second argument of $\tau$.

Thus, by \cref{thm:h-set-refrel-in-paths-sets} we have that $\susp(A)$ is a set and that $P(x,y) \eqvsym (\id{x}{y})$ for all $x,y:\susp(A)$.
Taking $x\defeq \north$ and $y\defeq \south$ yields $\eqv{A}{(\id[\susp(A)]\north\south)}$ as desired.
\end{proof}
